%%%%%%%%%%%%%%%%%
% This is an sample CV template created using altacv.cls
% (v1.7.2, 28 Aug 2024) written by LianTze Lim (liantze@gmail.com), based on the
% CV created by BusinessInsider at http://www.businessinsider.my/a-sample-resume-for-marissa-mayer-2016-7/?r=US&IR=T
%
%% It may be distributed and/or modified under the
%% conditions of the LaTeX Project Public License, either version 1.3
%% of this license or (at your option) any later version.
%% The latest version of this license is in
%%    http://www.latex-project.org/lppl.txt
%% and version 1.3 or later is part of all distributions of LaTeX
%% version 2003/12/01 or later.
%%%%%%%%%%%%%%%%

%% Use the "normalphoto" option if you want a normal photo instead of cropped to a circle
% \documentclass[10pt,a4paper,normalphoto]{altacv}

\documentclass[10pt,a4paper,ragged2e,withhyper]{altacv}
%% AltaCV uses the fontawesome5 and simpleicons packages.
%% See http://texdoc.net/pkg/fontawesome5 and http://texdoc.net/pkg/simpleicons for full list of symbols.

% Change the page layout if you need to
\geometry{left=1.25cm,right=1.25cm,top=1.5cm,bottom=1.5cm,columnsep=1.2cm}

% The paracol package lets you typeset columns of text in parallel
\usepackage{paracol}


% Change the font if you want to, depending on whether
% you're using pdflatex or xelatex/lualatex
% WHEN COMPILING WITH XELATEX PLEASE USE
% xelatex -shell-escape -output-driver="xdvipdfmx -z 0" mmayer.tex
\iftutex
  % If using xelatex or lualatex:
  \setmainfont{Lato}
\else
  % If using pdflatex:
  \usepackage[default]{lato}
\fi

% Change the colours if you want to
\definecolor{VividPurple}{HTML}{28518f}
\definecolor{SlateGrey}{HTML}{2E2E2E}
\definecolor{LightGrey}{HTML}{666666}
% \colorlet{name}{black}
% \colorlet{tagline}{PastelRed}
\colorlet{heading}{VividPurple}
\colorlet{headingrule}{VividPurple}
% \colorlet{subheading}{PastelRed}
\colorlet{accent}{VividPurple}
\colorlet{emphasis}{SlateGrey}
\colorlet{body}{LightGrey}

% Change some fonts, if necessary
% \renewcommand{\namefont}{\Huge\rmfamily\bfseries}
% \renewcommand{\personalinfofont}{\footnotesize}
% \renewcommand{\cvsectionfont}{\LARGE\rmfamily\bfseries}
% \renewcommand{\cvsubsectionfont}{\large\bfseries}

% Change the bullets for itemize and rating marker
% for \cvskill if you want to
\renewcommand{\cvItemMarker}{{\small\textbullet}}
\renewcommand{\cvRatingMarker}{\faCircle}
% ...and the markers for the date/location for \cvevent
% \renewcommand{\cvDateMarker}{\faCalendar*[regular]}
% \renewcommand{\cvLocationMarker}{\faMapMarker*}


% If your CV/résumé is in a language other than English,
% then you probably want to change these so that when you
% copy-paste from the PDF or run pdftotext, the location
% and date marker icons for \cvevent will paste as correct
% translations. For example Spanish:
% \renewcommand{\locationname}{Ubicación}
% \renewcommand{\datename}{Fecha}


%% Use (and optionally edit if necessary) this .tex if you
%% want to use an author-year reference style like APA(6)
%% for your publication list
% \input{pubs-authoryear.tex}

%% Use (and optionally edit if necessary) this .tex if you
%% want an originally numerical reference style like IEEE
%% for your publication list
\input{pubs-num.tex}

%% sample.bib contains your publications
\addbibresource{sample.bib}

\begin{document}
\name{Edgar Harutyunyan}
%\tagline{Business Woman \& Proud Geek}
% Cropped to square from https://en.wikipedia.org/wiki/Marissa_Mayer#/media/File:Marissa_Mayer_May_2014_(cropped).jpg, CC-BY 2.0
%% You can add multiple photos on the left or right
%\photoR{4cm}{edgar_photo}
% \photoL{2cm}{Yacht_High,Suitcase_High}
\personalinfo{%
  % Not all of these are required!
  % You can add your own with \printinfo{symbol}{detail}
  \email{harutyunianedgar@gmail.com}
  \phone{+49 (1523) 14-18-757}
  \location{Leipzig, Saxony}
  \mailaddress{Mainzer Str. 2, 04109 Germany}
  %\homepage{marissamayr.tumblr.com}
  % \twitter{@marissamayer}
  %\xtwitter{@marissamayer}
  \linkedin{edgar-harutyunian}
   \github{edhar98} % I'm just making this up though.
%   \orcid{0000-0000-0000-0000} % Obviously making this up too.
  %% You can add your own arbitrary detail with
  %% \printinfo{symbol}{detail}[optional hyperlink prefix]
  % \printinfo{\faPaw}{Hey ho!}
  %% Or you can declare your own field with
  %% \NewInfoFiled{fieldname}{symbol}[optional hyperlink prefix] and use it:
  % \NewInfoField{gitlab}{\faGitlab}[https://gitlab.com/]
  % \gitlab{your_id}
	%%
  %% For services and platforms like Mastodon where there isn't a
  %% straightforward relation between the user ID/nickname and the hyperlink,
  %% you can use \printinfo directly e.g.
  % \printinfo{\faMastodon}{@username@instace}[https://instance.url/@username]
  %% But if you absolutely want to create new dedicated info fields for
  %% such platforms, then use \NewInfoField* with a star:
  % \NewInfoField*{mastodon}{\faMastodon}
  %% then you can use \mastodon, with TWO arguments where the 2nd argument is
  %% the full hyperlink.
  % \mastodon{@username@instance}{https://instance.url/@username}
}
\makecvheader

%% Depending on your tastes, you may want to make fonts of itemize environments slightly smaller
\AtBeginEnvironment{itemize}{\small}

%% Set the left/right column width ratio to 6:4.
\columnratio{0.6}

% Start a 2-column paracol. Both the left and right columns will automatically
% break across pages if things get too long.
\begin{paracol}{2}
\cvsection{Education}

\cvevent{Bachelor of Science in International Physics Studies Program (Honours)}{\href{https://www.uni-leipzig.de/}{Leipzig University}}{Oct 2022 -- Sep 2026}{Leipzig, Germany}

\divider

\cvevent{Bachelor of Informatics in Information Technologies}{\href{https://www.polytech.am/en}{National Polytechnic University of Armenia}}{Sep 2016 -- Jul 2020}{Gyumri, Armenia}

\divider

\cvevent{High School Diploma}{Gyumri Academic College}{Sep 2013 - Jul 2016}{Gyumri, Armenia}

\cvsection{Experience}

\cvevent{Student Assistant at Faculty of Physics and Earth System Sciences}{\href{https://www.physes.uni-leipzig.de/}{Leipzig University}}{Oct 2024 -- Current}{Leipzig, Germany}

\divider

\cvevent{Software Engineer}{\href{https://ggg.instigatedesign.com/}{Instigate Design CJSC}}{Oct 2021 -- Jan 2025}{Yerevan, Armenia}

\textbf{\small{Remote since Feb 2023}}
\begin{itemize}
\item \textbf{Heterogeneous Computing and AI Applications}: Worked on a text semantic analysis project by deploying a classification algorithm on a hybrid system of FPGA accelerators and multi-core processors, achieving high-performance computation.
\item \textbf{Chip-to-Chip Data Transfer}: Contributed to designing and implementing a high-throughput system for transferring raw camera image data between processing modules in a satellite system.
\item \textbf{Satellite Image Analysis}: Developed an FPGA-based Ultra-High-Resolution (UHR) algorithm for real-time satellite image processing, facilitating advanced data analysis capabilities.
\item \textbf{Thread Library and OS Kernel Development}: Enhanced a multi-threading library for an OS kernel project by adding assembly support for RISC-V architecture.
\item \textbf{Custom Processor Design}: Designed and implemented custom CPU architecture, including stack-based and Von Neumann designs, with custom instruction set, debugging, and FPGA-based deployment.
\end{itemize}

\newpage
\cvsection{Training}

\cvevent{Participation in FPGA Conference 2024}{\href{https://www.fpga-conference.eu/}{ELEKTRONIKPRAXIS and the FPGA training center PLC2}}{02-04 Jul 2024}{Munich, Germany}

\divider

\cvevent{Machine Learning}{\href{https://bootcamps.aca.am/eng}{Armenian Code Academy}}{Dec 2020 -- April 2021}{Yerevan, Armenia}

\divider

\cvevent{Quality Engineering}{\href{https://www.bdg.am}{Business Development Group}}{Dec 2020 -- Feb 2021}{Yerevan, Armenia}

% \divider

% \cvevent{Product Engineer}{Google}{23 June 1999 -- 2001}{Palo Alto, CA}

% \begin{itemize}
% \item Joined the company as employe \#20 and female employee \#1
% \item Developed targeted advertisement in order to use user's search queries and show them related ads
% \end{itemize}

% use ONLY \newpage if you want to force a page break for
% ONLY the currentc column

%% Switch to the right column. This will now automatically move to the second
%% page if the content is too long.

\switchcolumn

\cvsection{Life Philosophy}
\begin{quote}
"You must accordingly be perfect, as your heavenly Father is perfect." - \\ Matthew 5:48

\end{quote}

\cvsection{Skills}
\cvskill{C/C++}{4.5}
\smallskip
{\LaTeXraggedright
\cvtag{STL}
\cvtag{OOP}
\cvtag{Gcc \& Gdb}
\cvtag{Parallel Programming (Threads, Processes)}
\par}
\divider
\cvskill{Python}{4}
\smallskip
{\LaTeXraggedright
\cvtag{OOP}
\cvtag{ML (Pandas, Scikit-learn, Etc.)}}
\cvtag{Data Analysis (Numpy, Matplotlib, Etc.)}


\divider
\cvskill{MATLAB}{3}
\divider
\cvskill{Low-Level Programming}{3}
\smallskip
{\LaTeXraggedright
\cvtag{x86/ARM}
\cvtag{MIPS}
\cvtag{RISC-V}
\par}
\divider\smallskip
\cvskill{Hardware Design}{4}
\smallskip
{\LaTeXraggedright
\cvtag{Verilog/VHDL}
\cvtag{Xilinx Vivado/Vitis}}

\divider\smallskip
% Don't overuse these \cvtag boxes — they're just eye-candies and not essential. If something doesn't fit on a single line, it probably works better as part of an itemized list (probably inlined itemized list), or just as a comma-separated list of strengths.

% The `ragged2e` document class option might cause automatic linebreaks between \cvtag to fail.
% Either remove the ragged2e option; or 
% add \LaTeXraggedright in the paragraph for these \cvtag
\cvskill{Development Tools \& Environments}{0}
\smallskip

{\LaTeXraggedright
\cvtag{Linux}
\cvtag{Bash}
\cvtag{GNU Make}
\cvtag{Vim}

\cvtag{Version control (Git, SVN)}
}

\divider\smallskip

\cvskill{Office \& Publishing Tools}{0}
\smallskip

{\LaTeXraggedright
\cvtag{MS Office Suite}
\cvtag{LaTeX}
\cvtag{OpenOffice / LibreOffice}

\cvtag{Google Workspace}
\cvtag{Adobe Acrobat}
}

\divider\smallskip

\cvskill{Team \& Agile Collaboration}{0}
\smallskip

{\LaTeXraggedright
\cvtag{TDD \& Code review}
\cvtag{Collaboration}
\cvtag{Communication}
\cvtag{Planning \& estimation}
\cvtag{Agile development practices}}

\newpage
\cvsection{Languages}

\cvlangskill{Armenian}{5} %% supports X.5 values.

\cvlangskill{English}{4}
% \divider
\cvlangskill{Russian}{4}
\cvlangskill{German (learning actively)}{1.5}
% \divider

\cvsection{Certification}
All certificates confirming the above data can be found by visiting the link below.

\textbf{\href{https://drive.google.com/drive/folders/1sEL77PSKGdF_0kmJpEL1gvbUvp6qNePs?usp=sharing}{https://shorturl.at/WTfj1}}

\end{paracol}

\cvsection{~}

\end{document}